\chapter{PTES Fase 7 - Reporting}

\section{Criterio de clasificacion}

Los hallazgos se clasifican con CVSS v3.1 cuando existe un impacto de
seguridad reproducible. Las observaciones de endurecimiento sin una cadena de
explotacion demostrada se mantienen como informativas y no reciben una
puntuacion artificial. Los controles efectivos se presentan de forma
separada para evidenciar que tambien fueron probados.

\section{Resumen ejecutivo de hallazgos}

\begin{center}
\captionof{table}{Resumen ejecutivo de hallazgos y severidad}\label{tab:hallazgos}
\end{center}
{\def\LTcaptype{none}
\begin{longtable}{@{}p{0.08\textwidth}p{0.26\textwidth}p{0.18\textwidth}p{0.10\textwidth}p{0.22\textwidth}@{}}
\toprule
\textbf{ID} & \textbf{Hallazgo} & \textbf{Vector CVSS v3.1} & \textbf{Puntaje} & \textbf{Clasificacion} \\
\midrule
\endfirsthead
\toprule
\textbf{ID} & \textbf{Hallazgo} & \textbf{Vector CVSS v3.1} & \textbf{Puntaje} & \textbf{Clasificacion} \\
\midrule
\endhead
F-01 &
El registro publico permite solicitar el rol \texttt{admin} y obtener acceso
administrativo sin autorizacion previa. &
\path{AV:N/AC:L/PR:N/UI:N/S:U/C:H/I:H/A:H} &
\textbf{9.8} &
\severity{4}{Critico} \\

F-02 &
El endpoint diagnostico \texttt{/api/db-test} esta publicado sin
autenticacion y puede revelar mensajes internos en la ruta de error. &
\path{AV:N/AC:L/PR:N/UI:N/S:U/C:L/I:N/A:N} &
\textbf{5.3} &
\severity{2}{Medio} \\

F-03 &
La carga de imagen confia en extension y MIME declarado; un archivo de texto
renombrado como JPG fue aceptado y servido sin
\texttt{X-Content-\allowbreak Type-Options}. &
\path{AV:N/AC:L/PR:L/UI:R/S:C/C:L/I:L/A:N} &
\textbf{6.1} &
\severity{2}{Medio} \\

F-04 &
La carga masiva convierte un valor numerico invalido a cero en lugar de
rechazar la fila, afectando la integridad de datos de inventario. &
\path{AV:N/AC:L/PR:H/UI:N/S:U/C:N/I:L/A:N} &
\textbf{2.7} &
\severity{1}{Bajo} \\

O-01 &
CORS responde con \texttt{Access-Control-Allow-\allowbreak Origin: *}, pero no habilita
credenciales cross-origin y los endpoints sensibles mantienen JWT. &
No aplica &
N/A &
\severity{0}{Observacion} \\

O-02 &
El frontend no define CSP ni \texttt{Permissions-Policy}. &
No aplica &
N/A &
\severity{0}{Endurecimiento} \\

O-03 &
El dominio local apunta al nodo Minikube; Nmap observo puertos de plataforma
fuera de la superficie web. &
No aplica &
N/A &
\severity{0}{Topologia} \\

O-04 &
El backend usa \texttt{DB\_USER: root}, lo que ampliaria el impacto de un
compromiso futuro del pod. &
No aplica &
N/A &
\severity{0}{Arquitectura} \\

O-05 &
No se deshabilita el montaje automatico del token de ServiceAccount y no
existe una politica de egreso por defecto. &
No aplica &
N/A &
\severity{0}{Kubernetes} \\
\bottomrule
\end{longtable}
}

\begin{tcolorbox}[criticalbox,title={F-01 - Escalacion de privilegios desde el registro}]
\textbf{Resultado.} Una solicitud anonima a
\texttt{POST /api/auth/register} acepto \texttt{rol: admin}; el login
posterior emitio un JWT administrativo y permitio consultar
\texttt{/api/usuarios}. La cuenta temporal fue desactivada al terminar.

\textbf{Impacto.} Compromiso de confidencialidad e integridad del inventario,
usuarios y auditoria, con capacidad administrativa sobre la aplicacion.

\textbf{Mitigacion.} Ignorar cualquier rol recibido desde el cliente,
asignar el rol minimo en el backend, cerrar el registro publico y reservar la
creacion de administradores para un flujo autenticado y auditado.
\end{tcolorbox}

\section{Controles de seguridad validados}

\begin{center}
\captionof{table}{Controles de seguridad validados}\label{tab:controles-validados}
\end{center}
{\def\LTcaptype{none}
\begin{longtable}{p{0.12\textwidth}p{0.38\textwidth}p{0.39\textwidth}}
\toprule
\textbf{ID} & \textbf{Control} & \textbf{Resultado observado} \\
\midrule
\endfirsthead
\toprule
\textbf{ID} & \textbf{Control} & \textbf{Resultado observado} \\
\midrule
\endhead
C-01 & Rate limiting del login &
Tres respuestas 401 y una cuarta 429 con
\texttt{RateLimit-Remaining: 0} y \texttt{Retry-After}. \\
C-02 & Validacion JWT &
401 sin token, 403 con token invalido o alterado y 200 con JWT valido. \\
C-03 & Revocacion por usuario inactivo &
El token previo fue rechazado con 403 despues de desactivar la cuenta. \\
C-04 & RBAC admin/almacenero &
El almacenero recibio 403 en usuarios, auditoria, eliminacion y restauracion. \\
C-05 & SQLi clasica en login &
La prueba no emitio token ni mostro errores SQL; se mantiene el uso de
consultas parametrizadas. \\
C-06 & Bcrypt &
Ocho registros mostraron longitud 60 y prefijo \texttt{\$2b\$}, sin exponer
hashes completos. \\
C-07 & Chatbot &
No revelo prompt ni secretos, trato la cadena SQL como texto y mantuvo el
alcance de datos de inventario. \\
C-08 & Contenedores y Kubernetes &
Backend y frontend operan sin root; el retest final mostro pods
\texttt{1/1 Running}, cero reinicios y health HTTP 200. \\
\bottomrule
\end{longtable}
}

\section{Recomendaciones priorizadas}

\begin{enumerate}
  \item \textbf{Inmediata:} corregir el registro para impedir que el cliente
  seleccione el rol y desactivar el registro publico de administradores.

  \item \textbf{Alta:} retirar o proteger \texttt{/api/db-test}; las rutas de
  error deben devolver mensajes genericos y registrar el detalle solo en logs
  internos.

  \item \textbf{Alta:} validar archivos por contenido real, extension, MIME y
  firma; almacenar fuera de rutas ejecutables, generar nombres seguros y
  devolver \texttt{X-Content-Type-Options: nosniff}.

  \item \textbf{Media:} validar tipos y rangos en carga masiva. Una fila con
  numero invalido debe rechazarse con un error identificable y no convertirse
  silenciosamente a cero.

  \item \textbf{Media:} limitar CORS a los origenes necesarios y definir una
  CSP compatible con Vite/React, junto con \texttt{Permissions-Policy}.

  \item \textbf{Continua:} mantener pruebas automatizadas para JWT, RBAC,
  rate limit, bcrypt y revocacion de sesion en cada entrega.

  \item \textbf{Alta:} crear un usuario MySQL exclusivo para la aplicacion
  con privilegios minimos y reservar \texttt{root} para administracion.

  \item \textbf{Media:} definir \texttt{automountServiceAccountToken: false}
  en workloads que no consumen la API de Kubernetes.

  \item \textbf{Media:} agregar una politica de egreso \texttt{default-deny}
  y permitir solo DNS, backend hacia MySQL y el destino externo necesario.
\end{enumerate}

\section{Plan de retest}

\begin{enumerate}
  \item Repetir el registro anonimo solicitando \texttt{rol: admin}; se espera
  rechazo o asignacion forzada del rol minimo.
  \item Confirmar que \texttt{/api/db-test} devuelve 404 o 401/403 y nunca
  propaga \texttt{error.message}.
  \item Repetir E05-14 con un marcador de texto renombrado; se espera 400/415.
  \item Repetir la carga masiva con \texttt{abc} en un campo numerico; se
  espera que la fila sea rechazada sin crear productos.
  \item Ejecutar nuevamente ZAP, cabeceras, health, pods y eventos.
  \item Repetir Nuclei con un conjunto web acotado que pueda finalizar; el
  barrido actual de 1\% no debe presentarse como cobertura completa.
  \item Ejecutar una matriz pod a pod solo tras aprobar el addendum de alcance;
  si el punto de entrada es \texttt{kubectl exec}, reportarla como simulacion
  de brecha y no como movimiento lateral explotado.
\end{enumerate}

\section{Conclusiones}

La practica completo el ciclo de contenerizacion, despliegue Kubernetes,
exposicion por Ingress y evaluacion PTES. El sistema fue comprometido a nivel
de aplicacion: una solicitud anonima creo una cuenta administradora, obtuvo un
JWT con rol \texttt{admin} y accedio a datos protegidos. Este escalamiento
vertical es una vulneracion real y constituye el hallazgo critico principal.

Los controles de bcrypt, JWT, revocacion por estado activo, RBAC, rate
limiting, consultas parametrizadas y aislamiento de contenedores respondieron
de acuerdo con las pruebas ejecutadas. No se demostro movimiento lateral
porque no se obtuvo ejecucion de codigo ni un punto de apoyo dentro de un pod.
El retest final confirmo que la limpieza no afecto la disponibilidad.

\section{Limitaciones}

\begin{itemize}
  \item El alcance se limito al laboratorio local autorizado y al dominio
  \texttt{conjunta3p.espe.edu.ec}.
  \item No se ejecutaron ataques destructivos, cracking de hashes, extraccion
  completa de la base ni pruebas contra la infraestructura de Groq.
  \item Burp Suite, Hydra y SQLMap no tienen evidencia de ejecucion. Sus
  casos se sustituyeron o descartaron con justificacion y no se les atribuyen
  resultados.
  \item No se ejecutaron movimiento lateral, acceso a kubelet/etcd/API Server,
  lectura de tokens de ServiceAccount, MySQL directo, RCE ni escape de
  contenedor. El borrador de ampliacion no autoriza pruebas.
  \item Nikto tuvo un timebox de 90 segundos; Nuclei cubrio solo 1\% de las
  solicitudes estimadas; ZAP completo un quick scan no autenticado. Estas
  limitaciones se complementaron con validacion manual.
  \item Los puntajes CVSS representan el impacto demostrado en este entorno y
  deben recalcularse si cambia la arquitectura de produccion.
\end{itemize}
